% ============================================================
% Lost in Transliteration: Why Romanization Breaks LLM Safety
% Draft: Introduction + Related Work (repositioned after the
% August 2026 Bangla safety benchmarks; see Related Work).
% Requires: \usepackage{natbib} (ACL style uses \citet / \citep)
%           \usepackage{xcolor}
% Placeholders marked \todo{...} must be filled from your results.
% Citations marked "VERIFY" in references.bib must be checked.
% ============================================================

\providecommand{\todo}[1]{\textcolor{red}{[#1]}}

\section{Introduction}

Large language models (LLMs) are aligned to refuse harmful requests through
supervised fine-tuning and reinforcement learning from human feedback
\citep{ouyang2022training, bai2022constitutional}, and are often further
protected by input--output guard models \citep{inan2023llamaguard,
zeng2024shieldgemma}. Yet safety alignment generalizes unevenly:
\citet{wei2023jailbroken} attribute many jailbreaks to \emph{mismatched
generalization}, where a model's capabilities reach input distributions that
its safety training does not. Language is a prominent example. Harmful
requests translated into low-resource languages are answered far more often
than their English originals \citep{yong2023lowresource,
deng2024multilingual, wang2024languages}.

A growing body of work shows that \emph{how} a language is written matters as
much as \emph{which} language it is. Transliterated Arabic and Arabizi bypass
safeguards that hold for Arabic script \citep{ghanim2024arabizi};
phonetically perturbed, code-mixed Hinglish and Banglish jailbreak frontier
models \citep{aswal2025cmprt}; and, for Bangla specifically, three recent
benchmarks document large safety gaps across scripts, registers and
spelling variants \citep{banglaveilguard2026, banglasafe2026,
alignedform2026}. Together, these studies establish \emph{that} romanized
and informal Bangla---commonly called \emph{Banglish}---weakens LLM safety.

What remains unclear is \emph{why}. Existing studies vary several factors at
once: a Banglish prompt differs from its English counterpart in language,
and from its Bangla-script counterpart in script, so a higher attack success
rate cannot be attributed to either. Explanations offered so far are
largely behavioral, pointing to tokenization \citep{aswal2025cmprt} or to
training-data imbalance \citep{alignedform2026}, while interpretability work
suggests that refusal is mediated by a direction that is shared across
languages \citep{arditi2024refusal, wang2025refusaluniversal}. Whether this
shared mechanism extends to a language written in a \emph{foreign script},
and where it breaks down, has not been tested. Finally, current defenses
for Bangla are prompt guards \citep{banglaveilguard2026}; their robustness to
an attacker who knows the defense \citep{nasr2025attacker} is unknown.

We address these questions with a controlled study. Our key design choice is
a $2 \times 2$ factorial crossing \textbf{language} (English, Bangla) with
\textbf{script} (Latin, Bengali): each seed request is written in English
(Latin script), English transliterated into Bengali script, Bangla (Bengali
script) and Banglish (Bangla in Latin script). Holding intent fixed, this
isolates the effect of script from the effect of language. Two further
versions---Banglish with real-world spelling noise and Banglish--English
code-mixing---extend the design to the forms users actually type. Because a
model that fails to refuse may simply have misunderstood the input, we score
every response on two axes, comprehension and harmfulness, and count only
on-topic harmful responses as successful attacks.

Our contributions are:
\begin{enumerate}
    \item \textbf{A script--language factorial benchmark.} \todo{N} native-authored
    seeds, each in six parallel versions, with benign controls for
    over-refusal and Bangladesh-specific harm categories. Using it, we show
    that \todo{e.g., ``script, not language, accounts for most of the safety
    gap: moving Bangla from Bengali to Latin script raises the attack
    success rate by X points, while moving English to Bengali script raises
    it by only Y''}.

    \item \textbf{Comprehension-controlled measurement.} Across \todo{M}
    proprietary and open-weight models and \todo{G} guard models, we separate
    genuine jailbreaks from misunderstanding, and quantify a dose--response
    relation between real-world spelling noise and attack success
    \todo{summarize}.

    \item \textbf{A mechanistic account.} Extending the cross-lingual
    refusal-direction analysis of \citet{wang2025refusaluniversal} to
    scripts, we find that \todo{e.g., ``harmful Banglish is represented
    close to its English counterpart in middle layers, yet projects X\%
    less onto the refusal direction''}, and that adding the direction
    restores refusals \todo{causal result}.

    \item \textbf{Defenses under adaptive attack.} We evaluate normalization
    before moderation, safety fine-tuning and activation steering, and
    show that \todo{e.g., ``existing Bangla prompt guards, including
    BanglaVeilGuard, can be evaded by adaptive spelling attacks, while
    ...''}. We also report results on the BanglaVeilGuard and BanglaSafe
    evaluation sets to test generalization.
\end{enumerate}

We release code openly and the benchmark under gated access, and disclosed
our findings to the affected model providers before publication. The
factorial design and analysis apply directly to other widely romanized
languages such as Hindi (Hinglish), Urdu and Arabic (Arabizi).

% ============================================================
\section{Related Work}

\paragraph{Jailbreak attacks on aligned LLMs.}
Manually crafted jailbreaks have been collected and characterized at scale
\citep{shen2024anything}. \citet{wei2023jailbroken} identified competing
objectives and mismatched generalization as two root causes of safety
failures; our study targets the latter. Automated methods include
gradient-based suffix optimization \citep{zou2023universal}, attacker-LLM
refinement such as PAIR \citep{chao2023jailbreaking} and TAP
\citep{mehrotra2024tree}, persuasion \citep{zeng2024johnny} and many-shot
prompting \citep{anil2024manyshot}. HarmBench \citep{mazeika2024harmbench}
and JailbreakBench \citep{chao2024jailbreakbench} standardize evaluation,
but are English-only.

\paragraph{Multilingual safety.}
Translating unsafe prompts into low-resource languages bypasses safeguards
\citep{yong2023lowresource}, with unsafe output increasing as language
resources decrease \citep{deng2024multilingual}. XSafety
\citep{wang2024languages} and further analyses \citep{shen2024language}
study multilingual safety in native scripts, and human-written multilingual
red-teaming data supports local harm mitigation
\citep{aakanksha2024multilingual}. Code-switching red-teaming
\citep{yoo2024codeswitching} shows that mixing languages within a prompt
elicits unsafe behavior.

\paragraph{Script, transliteration and Bangla safety.}
\citet{ghanim2024arabizi} showed that Arabic written in transliteration or
Arabizi elicits unsafe responses that Arabic script does not.
\citet{aswal2025cmprt} combined code-mixing with phonetic misspellings of
sensitive words in Hinglish, and related work reports the same effect for
Banglish, explaining it by tokenization. Three recent Bangla benchmarks are
closest to our work. BanglaVeilGuard \citep{banglaveilguard2026} covers six
written forms of Bangla, including romanized, Banglish, code-mixed, noisy
and dialectal text, and proposes a lightweight prompt guard. BanglaSafe
\citep{banglasafe2026} anchors 17 harm categories in Bangladeshi statutes and
documented cases and finds that register (formal versus colloquial) shifts
attack success more than language does. \citet{alignedform2026} separate
comprehension from containment for Bangla derogatory speech and report a
``romanization paradox'' in which a model understands Banglish better than
Bangla script. We build on these findings but differ in three ways: our
factorial design separates script from language for the same intent;
we explain the gap mechanistically rather than behaviorally; and we test
defenses, including existing Bangla guards, against adaptive attackers.
\todo{Read the full papers and refine this paragraph; check whether any of
them uses parallel versions of the same seed.}

\paragraph{Romanization inside LLMs.}
LLMs appear to process some non-Latin-script languages partly through
romanized representations: romanized input can improve cross-lingual
transfer \citep{husain2024romansetu}, and intermediate layers show latent
romanization \citep{saji2025romanlens}. Resources such as the Dakshina
dataset \citep{roark2020dakshina} and IndicXlit \citep{madhani2023aksharantar}
support transliteration between scripts for South Asian languages. These
findings motivate our hypothesis that Banglish is understood through the
same pathways as English text while escaping the safety behavior learned
for it.

\paragraph{Mechanisms of refusal and defenses.}
Refusal in chat models is mediated by a single residual-stream direction
\citep{arditi2024refusal}, and this direction is shared across
safety-aligned languages; cross-lingual jailbreaks have been linked to weaker
separation of harmful and harmless inputs along it
\citep{wang2025refusaluniversal}. Safety alignment is also shallow
\citep{qi2025safety} and fragile under fine-tuning \citep{qi2024finetuning}.
Defenses include guard models \citep{inan2023llamaguard, zeng2024shieldgemma},
input smoothing \citep{robey2023smoothllm} and classifier-based safeguards
\citep{sharma2025constitutional}; \citet{nasr2025attacker} show that many
defenses fail against adaptive attackers. We use the refusal direction both
to explain the Banglish gap and as a steering defense, and evaluate all
defenses adaptively.

\paragraph{Bangla NLP.}
Despite its large speaker population, Bangla is low-resource in NLP
\citep{joshi2020state}; pretrained models such as BanglaBERT
\citep{bhattacharjee2022banglabert} and recent Bangla LLMs evaluated in
\citet{banglaveilguard2026} have narrowed but not closed this gap.
